\documentclass{article}

% Official ICLR 2026 style and bibliography format from the prior COMPOSE submission.
\usepackage{iclr2026_conference,times}
\usepackage[utf8]{inputenc}
\usepackage[T1]{fontenc}
\usepackage{microtype}
\usepackage{amsmath,amssymb,mathtools,bm}
\usepackage{amsthm}
\usepackage{booktabs,tabularx,array}
\usepackage{xspace}
\usepackage{float}
\usepackage{graphicx}
\usepackage{xcolor}
\usepackage{enumitem}
\usepackage{pifont}
\usepackage{url}
\usepackage{hyperref}
\usepackage[capitalize,noabbrev]{cleveref}

\hypersetup{colorlinks=true,citecolor=blue!55!black,linkcolor=blue!55!black,urlcolor=blue!55!black}
% Centralized placeholders for the final confirmatory evaluation.
% Replace XXX only after the corresponding artifact is frozen and audited.
\newcommand{\XXX}{\textbf{XXX}}

% Historical de novo placeholders. The reversible null source means endpoint and all-step
% non-null rates must be measured; the closure theorem does not authorize 100% empirical cells.
\newcommand{\ComposeValidity}{\XXX}
\newcommand{\ComposeTrajectoryValidity}{\XXX}
\newcommand{\ComposeTrajectoryConnectivity}{\XXX}
\newcommand{\ComposeUniqueness}{\XXX}
\newcommand{\ComposeNovelty}{\XXX}
\newcommand{\ComposeDescriptorW}{\XXX}
\newcommand{\ComposeDiversity}{\XXX}
\newcommand{\ComposeCoverage}{\XXX}
\newcommand{\ComposeMeanEvents}{\XXX}
\newcommand{\ComposeThroughput}{\XXX}
% FCD / GuacaMol KL are deferred to the final GuacaMol table; NOT reported in the
% CNOF control paper (we do not stage an FCD arms race).
\newcommand{\ComposeFCD}{\XXX}
\newcommand{\ComposeKL}{\XXX}

% Structural distribution metrics.
\newcommand{\ComposeSizeTV}{\XXX}
\newcommand{\ComposeCycleRankTV}{\XXX}
\newcommand{\ComposeRingSizeTV}{\XXX}
\newcommand{\ComposeElementTV}{\XXX}
\newcommand{\ComposeBondTV}{\XXX}

% Operator and coupling diagnostics.
\newcommand{\ComposeGraftFraction}{\XXX}
\newcommand{\ComposeGrowFraction}{\XXX}
\newcommand{\ComposeShrinkFraction}{\XXX}
\newcommand{\ComposeRetypeFraction}{\XXX}
\newcommand{\ComposeRingFraction}{\XXX}
\newcommand{\ComposePathLength}{\XXX}
\newcommand{\ComposeCatalogCoverage}{\XXX}

% Ablations.
\newcommand{\NullPriorFCD}{\XXX}
\newcommand{\PrimitiveTreeFCD}{\XXX}
\newcommand{\GraftTreeFCD}{\XXX}
\newcommand{\NoMacroFCD}{\XXX}
\newcommand{\NoAliasFCD}{\XXX}
\newcommand{\NoStratificationFCD}{\XXX}

% Superseded Base-B/CNOF control results are quarantined and cannot populate this manuscript.
\newcommand{\ComposeExactnessLinf}{\XXX}
\newcommand{\ComposeScaffoldSatisfaction}{\XXX}
\newcommand{\ComposeSubstructureSatisfaction}{\XXX}
\newcommand{\ComposeScaffoldRetention}{\XXX}
\newcommand{\ComposeBaselineScaffoldEasy}{\XXX}
\newcommand{\ComposeBaselineScaffoldHard}{\XXX}
\newcommand{\ComposeOracleEfficiency}{\XXX}
\newcommand{\ComposeBaselineOracleEfficiency}{\XXX}
\newcommand{\ComposeBaselineOracleWastedHard}{\XXX}
\newcommand{\ComposeConstrainedQEDFrom}{\XXX}
\newcommand{\ComposeConstrainedQEDTo}{\XXX}
\newcommand{\ComposeConstrainedQEDGain}{\XXX}
\newcommand{\ComposeConstrainedLeadsImproved}{\XXX}
\newcommand{\ComposeAnytimeCalls}{\XXX}
\newcommand{\ComposePropertyGain}{\XXX}
\newcommand{\ComposeConditionalSuccess}{\XXX}

% Reproducibility.
% RETRACTED 2026-07-28: "6.3M" had no source. No script counts parameters, no checkpoint metadata
% field records one, and no diagnostics/results JSON contains a parameter-count key; both
% EXPERIMENT_CHECKLISTs still list "Parameter count" as unchecked. The count is also
% configuration-dependent (cycle heads exist only under enable_cycle_ops; atom heads are 15-wide
% under ORGANIC_VOCABULARY vs 4-wide under CNOF), so a base-B figure would not be RingCore-V1's.
\newcommand{\ComposeParameters}{\XXX}
\newcommand{\ComposeTrainHours}{\XXX}
\newcommand{\ComposePeakMemory}{\XXX}
% RETRACTED 2026-07-28: the "~560 oracle calls" figure for the oracle-guided controller had no
% artifact anywhere in repo history (it traces to prose in docs/CONDITIONAL_RESULTS_AND_SCOPING.md).
% Committed per-lead means across the unconstrained panels are 458-587; pick and freeze one.
\newcommand{\ComposeOracleHungryCalls}{\XXX}
% RETRACTED 2026-07-28: no artifact backs any of the developmental compiler/training figures
% (99,738 programs; 49,869/50,000 targets; RGM loss 84.09->28.14; "3,000-update checkpoint";
% 90.2% teacher family accuracy). The 90.2% value exists in
% results/tree_fcd_transfer_stage1_factorized_v1/metrics.json but for a run with 8,000 steps,
% train_size 8,000, and loss 34.98->-1.93, i.e. not the run the sentence describes.
\newcommand{\DevCompiledPrograms}{\XXX}
\newcommand{\DevCompiledTargets}{\XXX}
\newcommand{\DevTargetTotal}{\XXX}
\newcommand{\DevUpdates}{\XXX}
\newcommand{\DevLossFrom}{\XXX}
\newcommand{\DevLossTo}{\XXX}
\newcommand{\DevFamilyAccuracy}{\XXX}

% RingCore-V1 production training schedule (owner-locked 2026-07-28; scheduler_config_hash 0b832985c65de1cc).
% These are FROZEN (not XXX): AdamW, cosine LR with linear warmup. schedule_steps is the cosine-decay horizon
% matched to the ~1-3k warm-start convergence scale (NOT the base-B de-novo 30000). The schedule-faithful
% check (1000 steps) and the eventual full run share this schedule; only training_steps differs.
\newcommand{\RingCoreOptimizer}{AdamW}
\newcommand{\RingCorePeakLR}{$3\times10^{-4}$}
\newcommand{\RingCoreWeightDecay}{$1\times10^{-5}$}
\newcommand{\RingCoreWarmupSteps}{500}
\newcommand{\RingCoreScheduleSteps}{3000}
\newcommand{\RingCoreMinLRFraction}{0.05}
\newcommand{\RingCoreSchedule}{cosine with linear warmup}
\newcommand{\RingCoreTrainingSteps}{3000}

% ============================================================================
% RGM REFRAME (2026-07-28) --- macros introduced by the RGM-first restructure.
% Every RESULT macro for the RingCore-V1 production model is \XXX by design: that model is
% training at the time of writing and none of its results exist yet. Do NOT fill any of them
% from a preflight artifact --- the bounded-run, schedule-check, continuation, and rollout-panel
% artifacts are tagged NON_SCIENTIFIC_PREFLIGHT and are excluded from the paper by contract.
% Blocks that carry REAL values are (a) frozen configuration facts (scope, hashes, schedule),
% and (b) values read from committed non-preflight artifacts; each is labelled and sourced below.
% ============================================================================

% --- Production state-space / data contract (FROZEN facts, not results) ---
\newcommand{\ScopeMaxAtoms}{40}
\newcommand{\ScopeClasses}{15}
\newcommand{\ScopeHash}{\texttt{3721d69851110fdd}}
\newcommand{\ContractFingerprint}{\texttt{3917a0bf4a3de9830722663f}}
\newcommand{\CapabilityHash}{\texttt{330473e319bfec19}}
\newcommand{\OperatorRegistryHash}{\texttt{9197401e8dc3a7ae}}
\newcommand{\RingCatalogHash}{\texttt{639ff6078c32d43c}}
\newcommand{\TeacherFilterHash}{\texttt{f8137995181eeadc}}
% Corpus census under the production scope (measured; scaled_edit_data_manifest_40 / SYSTEM_CONTRACT 0).
\newcommand{\CorpusRetained}{466{,}483}
\newcommand{\CorpusTotal}{500{,}000}
\newcommand{\CorpusRetainedPct}{93.3\%}
% Base-B de novo pretraining subset (CNOF-neutral), SYSTEM_CONTRACT 0.
\newcommand{\BaseCnofSubset}{50{,}000}
\newcommand{\BaseCnofEligible}{\XXX}  % unused in prose; two repo sources disagree (225,149 vs 239,867) -- re-measure before use
% Teacher-representability filter removal rate (measured on 400 held-out broad-organic molecules).
\newcommand{\TeacherFilterRemoval}{\XXX}

% --- Fig. 2 / Table 1: unconditional (de novo) generative modeling. ALL PENDING. ---
\newcommand{\DenovoValid}{\XXX}
\newcommand{\DenovoConnected}{\XXX}
\newcommand{\DenovoUnique}{\XXX}
\newcommand{\DenovoNovel}{\XXX}
\newcommand{\DenovoFCD}{\XXX}
\newcommand{\DenovoPrecision}{\XXX}
\newcommand{\DenovoRecall}{\XXX}
\newcommand{\DenovoCoverage}{\XXX}
\newcommand{\DenovoScaffoldNovelty}{\XXX}
\newcommand{\DenovoIntDiv}{\XXX}
\newcommand{\DenovoSeeds}{\XXX}
\newcommand{\DenovoSampleCount}{\XXX}
% Non-memorization.
\newcommand{\DenovoNearestNeighbour}{\XXX}
\newcommand{\DenovoDuplicateGrowth}{\XXX}
\newcommand{\DenovoScaffoldOverlap}{\XXX}
% Property / size marginals (Wasserstein or TV against the held-out reference).
\newcommand{\DenovoMWDist}{\XXX}
\newcommand{\DenovoLogPDist}{\XXX}
\newcommand{\DenovoTPSADist}{\XXX}
\newcommand{\DenovoRotBondDist}{\XXX}
\newcommand{\DenovoChargeDist}{\XXX}
\newcommand{\DenovoHeavyAtomDist}{\XXX}
% Ring taxonomy marginals (the base-B ring defect is the reason RingCore exists).
\newcommand{\DenovoSmallRing}{\XXX}
\newcommand{\DenovoSmallRingRef}{\XXX}
\newcommand{\DenovoAromatic}{\XXX}
\newcommand{\DenovoAromaticRef}{\XXX}
\newcommand{\DenovoFused}{\XXX}
\newcommand{\DenovoSpiro}{\XXX}
\newcommand{\DenovoBridged}{\XXX}
\newcommand{\DenovoMacrocycle}{\XXX}

% --- Fig. 2f/2g/2h: learned transport vs. uniform legal rewrites; ring support; scaling. ---
\newcommand{\LearnedPathLength}{\XXX}
\newcommand{\UniformPathLength}{\XXX}
\newcommand{\LearnedNetDisplacement}{\XXX}
\newcommand{\UniformNetDisplacement}{\XXX}
\newcommand{\LearnedReversalRate}{\XXX}
\newcommand{\UniformReversalRate}{\XXX}
\newcommand{\LearnedEndpointFidelity}{\XXX}
\newcommand{\UniformEndpointFidelity}{\XXX}
\newcommand{\HeldOutRingTopology}{\XXX}
\newcommand{\RingCompositionalOnly}{\XXX}
\newcommand{\RingCatalogOnly}{\XXX}
\newcommand{\RingCompositionalPlusMacros}{\XXX}
\newcommand{\CanonicalBranching}{\XXX}
\newcommand{\ScalingData}{\XXX}
\newcommand{\ScalingBudget}{\XXX}
\newcommand{\ScalingModelSize}{\XXX}

% --- Fig. 3: exact finite-horizon control on enumerable molecular spaces. ---
% \ComposeExactnessLinf (above) is the COMMITTED continuous-time h-transform anchor.
\newcommand{\ComposeExactnessStates}{\XXX}
\newcommand{\ExactChainStates}{\XXX}
\newcommand{\ExactBudget}{\XXX}
\newcommand{\ExactTiltTV}{\XXX}
\newcommand{\ExactSwitchTV}{\XXX}
\newcommand{\ExactEarlyStopTV}{\XXX}
\newcommand{\ExactRareMotifTV}{\XXX}
\newcommand{\ExactIntervalTV}{\XXX}
\newcommand{\ExactMultiConstraintTV}{\XXX}
\newcommand{\ExactBoltzmannTV}{\XXX}
\newcommand{\ExactRegionTV}{\XXX}
\newcommand{\GreedyTV}{\XXX}
\newcommand{\LocalBoltzmannTV}{\XXX}
\newcommand{\RerankTV}{\XXX}
\newcommand{\MogStyleTV}{\XXX}
\newcommand{\ParticleTV}{\XXX}
\newcommand{\ExactHParticleFactor}{\XXX}
\newcommand{\SupportViolations}{\XXX}
\newcommand{\LearnedValueCalibration}{\XXX}
\newcommand{\BackwardResidual}{\XXX}
\newcommand{\RestartCost}{\XXX}
\newcommand{\ContinuationCost}{\XXX}

% --- Fig. 4 / Tables 2--3: source-conditioned multi-objective optimization. ALL PENDING. ---
\newcommand{\EditBudget}{\XXX}
\newcommand{\OracleBudget}{\XXX}
\newcommand{\ReferencePoint}{\XXX}
\newcommand{\EditHV}{\XXX}
\newcommand{\EditHVAUC}{\XXX}
\newcommand{\EditIGDPlus}{\XXX}
\newcommand{\EditFeasibleND}{\XXX}
\newcommand{\EditSourceSimilarity}{\XXX}
\newcommand{\EditOracleEfficiency}{\XXX}
\newcommand{\EditPathValidity}{\XXX}
\newcommand{\AdaptationRegret}{\XXX}
\newcommand{\SwitchCost}{\XXX}
\newcommand{\FanBranches}{\XXX}
\newcommand{\PathwiseFeasibleFraction}{\XXX}
\newcommand{\TerminalFilterFeasibleFraction}{\XXX}
\newcommand{\HeldOutOracleGap}{\XXX}
% Citations still to be verified against primary sources before submission.
\newcommand{\XXXcite}{\textbf{[CITE-XXX]}}

% ============================================================================
% COMMITTED, NON-PREFLIGHT ARTIFACTS newly cited by the RGM reframe.
% Each value below was read out of a committed JSON under diagnostics/ or results/.
% ============================================================================

% Benchmark-lead scope coverage -- diagnostics/composition/benchmark_lead_scope_coverage.json
\newcommand{\LeadCoverageBroad}{800/800}
\newcommand{\LeadCoverageCnof}{33.0\%}
% Broad corpus census under the production scope -- diagnostics/composition/broad_characterization_40.json
\newcommand{\MMPPairs}{216{,}183}
\newcommand{\MMPPoolRecords}{363{,}456}
\newcommand{\ReplayGate}{10{,}096/10{,}096}

% Conjunctive design specs -- diagnostics/conditional_smc/conjunction_funnel.json
\newcommand{\ConjunctionLeads}{4}
\newcommand{\ConjunctionClauses}{7}
\newcommand{\ConjunctionFullSpec}{8.7\%}
\newcommand{\ConjunctionDrawsPerHit}{11.4}

% Doob-vs-bootstrap particle convergence on the enumerated slice --
% diagnostics/exactness/doob_guidance_ground_truth.json (median TV over 12 replicates).
\newcommand{\ExactHParticles}{1{,}000}
\newcommand{\ExactHTV}{0.060}
\newcommand{\BootstrapTV}{0.116}

% Exact reachability / transport audits (results/).
\newcommand{\ReachabilityExact}{496/496}
\newcommand{\ReachabilityTargets}{124}
\newcommand{\RingTransportAudit}{200/200}
\newcommand{\TeacherFiberCoverage}{903/903}

% Objective-barrier diagnostics on exactly enumerated rewrite graphs -- diagnostics/reachability/.
\newcommand{\ReachStatesCapFour}{739}
\newcommand{\ReachEdgesCapFour}{9{,}870}
\newcommand{\ReachStatesCapFive}{5{,}454}
\newcommand{\BarrierSuboptTwoObj}{18.6\%}
\newcommand{\BarrierSuboptThreeObj}{27.4\%}
\newcommand{\BarrierSuboptWorstTwoObj}{33.7\%}
\newcommand{\BarrierSuboptCapFive}{23.9\%}
\newcommand{\BarrierSuboptWorstCapFive}{91.1\%}
\newcommand{\BarrierLookaheadTwo}{16.0\%}
\newcommand{\BarrierLookaheadThree}{8.9\%}
\newcommand{\BarrierInterior}{17.0\%}
\newcommand{\BarrierBoundary}{18.2\%}
\newcommand{\BarrierObjectivePairs}{7}
% The matched-oracle-budget frontier-RECOVERY comparison is deliberately absent: an earlier
% "detours lose at matched budget" reading was retracted as tuning-dependent. Efficiency of
% recovery is stated as undecided, never as a result.

% ============================================================================
% OPERATOR-SYSTEM SUPPORT CHECKS (executor-verified properties of the rewrite system, NOT model
% results) sourced from diagnostics/production_preflight/. These files carry no
% NON_SCIENTIFIC_PREFLIGHT tag, unlike the bounded-run and rollout-panel artifacts in the same
% directory. OWNER DECISION REQUIRED: if preflight-directory artifacts are ruled paper-ineligible
% in full, set every macro in this block to \XXX -- nothing else in main.tex needs to change.
% ============================================================================
\newcommand{\RingTopologyClasses}{14/14}          % ring_topology_capability.json
\newcommand{\RingRoundTripMolecules}{400}          % ring_compositional_support.json
\newcommand{\RingRoundTripEdges}{1{,}346}          % ring_compositional_support.json
\newcommand{\RingPathCostSystems}{800}             % ring_path_cost.json
\newcommand{\RingRestructureMedian}{3}             % ring_path_cost.json
\newcommand{\RingRestructureMax}{7}                % ring_path_cost.json
\newcommand{\RingRestructureOverBudget}{0\%}       % ring_path_cost.json
\newcommand{\RingIntroductionMedian}{20}           % ring_path_cost.json
\newcommand{\RingIntroductionOverBudget}{73.5\%}   % ring_path_cost.json
\newcommand{\RingEditBudget}{16}                   % ring_path_cost.json (production sampler horizon)


% ============================================================================
% DATA-STARVATION RECORD (diagnostics/coherence/program_state.json, 2026-07-28).
% These are counts of the training CONFIGURATION, not model results, and they are what bars every
% RingCore checkpoint from scientific status. Reported in app:status.
% ============================================================================
\newcommand{\StarvedRecords}{6{,}922}
\newcommand{\StarvedCorruptionSources}{200}
\newcommand{\StarvedMMPTraces}{5{,}000}
\newcommand{\StarvedPasses}{28}

% --- Experiments C (trans-dimensional) and D (topological adaptation): ALL PENDING. ---
\newcommand{\CardSmallerTarget}{\XXX}
\newcommand{\CardLargerTarget}{\XXX}
\newcommand{\CardGrowThenDelete}{\XXX}
\newcommand{\CardHeldOutBins}{\XXX}
\newcommand{\CardNoInsert}{\XXX}
\newcommand{\CardNoDelete}{\XXX}
\newcommand{\TopoCycleCreate}{\XXX}
\newcommand{\TopoCycleOpen}{\XXX}
\newcommand{\TopoRestructure}{\XXX}
\newcommand{\TopoHeldOut}{\XXX}
% --- Experiment E (quotient invariance): ALL PENDING. ---
\newcommand{\QuotientAggregateMatch}{\XXX}
\newcommand{\QuotientReencodeTV}{\XXX}
\newcommand{\QuotientControlledTV}{\XXX}
\newcommand{\QuotientMarkLevelTV}{\XXX}


\newcommand{\method}{\textsc{Compose}\xspace}
\newcommand{\rgm}{\textsc{RGM}\xspace}
\newcommand{\M}{\mathcal{M}_{\Gamma}}
\newcommand{\Rset}{\mathcal{R}}
\newcommand{\A}{\mathcal{A}}
\newcommand{\Sset}{\mathcal{S}}
\newcommand{\Z}{\mathcal{Z}}
\newcommand{\pdata}{p_{\mathrm{data}}}
\newcommand{\pzero}{\pi_0}
\newcommand{\T}{\mathsf{T}}
\newcommand{\legal}{\mathsf{Legal}}
\newcommand{\Match}{\mathsf{Match}}
\newcommand{\cmark}{\ding{51}}
\newcommand{\xmark}{\ding{55}}
\newcommand{\partialmark}{\ensuremath{\triangle}}
\newcommand{\E}{\mathbb{E}}
\newcommand{\R}{\mathbb{R}}
\newcommand{\one}{\mathbf{1}}
\newcolumntype{Y}{>{\raggedright\arraybackslash}X}

\newtheorem{theorem}{Theorem}
\newtheorem{proposition}[theorem]{Proposition}
\newtheorem{lemma}[theorem]{Lemma}
\newtheorem{corollary}[theorem]{Corollary}
\theoremstyle{definition}
\newtheorem{definition}[theorem]{Definition}
\newtheorem{assumption}[theorem]{Assumption}
\theoremstyle{remark}
\newtheorem{remark}[theorem]{Remark}
\crefname{assumption}{assumption}{assumptions}
\Crefname{assumption}{Assumption}{Assumptions}

\title{\method: Generator-Matched Stochastic Rewriting for Valid Molecular Generation and Editing}
\author{Anonymous Authors}
\date{}
% \iclrfinalcopy

\begin{document}
\maketitle

\begin{abstract}
Most molecular graph generators constrain an endpoint, not the states through which it is produced. This is a consequential distinction when generation is stopped early, guided by a property oracle, or initialized from an existing molecule. We introduce \emph{Rewrite Generator Matching} (\rgm), in which the infinitesimal events of a generative process are executable molecular graph rewrites. Its molecular instantiation, \method, is a continuous-time Markov chain (CTMC) on complete, connected, valence-admissible molecular graphs. Hard application conditions give illegal moves zero generator support; a neural model learns only the context- and time-dependent rates of legal insertion, deletion, relabeling, bond, subtree-Graft, and coordinated ring events. Training couples a directly samplable valid-graph prior to data and compiles executable source-to-target programs. These programs induce tractable conditional jump generators; Generator Matching marginalizes the endpoint and program, leaving a target-free generator at inference. Sampling is therefore ancestral CTMC simulation---not beam search, target-conditioned planning, or terminal repair. On GuacaMol, \method obtains FCD \ComposeFCD, validity \ComposeValidity, uniqueness \ComposeUniqueness, and all-step validity/connectivity \ComposeTrajectoryValidity/\ComposeTrajectoryConnectivity. All pending confirmatory measurements are deliberately marked \textbf{XXX}.
\end{abstract}

\section{Introduction}

A molecular generator defines more than an endpoint distribution: it defines how probability moves through chemical state space. Nevertheless, most molecular objectives supervise atoms, bonds, coordinates, or denoised graph tensors and ask whether the final decode sanitizes. Intermediate states may be disconnected, over-valent, or not molecular graphs at all. That can be acceptable when the path is purely latent numerical bookkeeping, but it is a structural limitation when generation is interrupted, scored online, initialized from a lead, restricted by a scaffold, or reused as an editor. In these settings, endpoint sanitization cannot make an undefined intermediate descriptor meaningful.

Graph diffusion and flow models have substantially advanced molecular distribution learning \citep{jo2022gdss,vignac2023digress,jo2023grum,xu2024disco,siraudin2024cometh,qin2025defog}. Flexible-cardinality methods additionally insert or delete atoms \citep{ninniri2025griddd,havasi2025editflows,franke2026morph}, and discrete bridges learn transformations between graph distributions \citep{kim2025ddsbm}. Their probability paths are primarily defined over node/edge variables, variable-length objects, or geometric coordinates. Separately, molecular graph grammars encode constraints in derivations \citep{kajino2019mhg,guo2022grammar}, while stochastic graph rewriting gives rule matches CTMC semantics under application conditions \citep{danos2015thermodynamic,behr2021rewriting}. We ask whether these ideas can be combined without reducing one to terminology: can the \emph{rewrite system itself} be the learned generative process, with chemistry defining which jumps exist and data determining their contextual firing rates?

We answer with \emph{Rewrite Generator Matching} (\rgm), instantiated for molecules as \method---\textbf{COM}positional \textbf{P}robabilistic \textbf{O}perator \textbf{S}ampling and \textbf{E}diting. A complete molecular graph is the state; a typed rule at a concrete match is a marked jump; and a neural network predicts its state- and time-dependent rate. Application conditions remove illegal actions before normalization. If symmetric matches or distinct rule lowerings reach the same chemical graph, their rates are summed in the quotient generator on molecules. Thus stochastic rewriting is operational rather than decorative: it specifies the state graph, legality, event identity, inverse operations, and ancestral simulator.

The learning obstacle is that a target molecule admits many valid edit programs and no closed-form posterior over them. We sample a source $G_0$ from a valid prior $\pzero$, a target $G_1$ from the data distribution, and a certified program connecting the pair. A progress CTMC on the program has exact conditional rates. Generator Matching \citep{holderrieth2025generator} regresses those rates and marginalizes the endpoint, source--target coupling, and latent program. The compiler is privileged supervision, not a decoder: inference receives only the current molecule and time.

The universal micro-operator set is deliberately small, but primitive edit paths can be poorly coordinated. A random carbon tree can reach a target by deleting to one atom and regrowing, yet this destroys useful topology and encourages an effectively constructive teacher. \method therefore adds two classes of verified derived rules. A \emph{subtree Graft} atomically exchanges a tree bridge and reconnects the detached component, making global topology revisable without a disconnected intermediate. \emph{Coordinated ring rewrites} introduce cycles and open ears as single semantic events and jointly restate common ring electronics. These macros do not replace the primitive support: they lower to validity-checked micro programs and have verified inverses.

Our contributions are:
\begin{itemize}[leftmargin=*]
    \item \textbf{Rewrite Generator Matching.} We define a jump-process Generator Matching construction whose mark space is the legal matches of a stochastic graph rewrite system, including the required quotient from syntactic marks to chemical successors.
    \item \textbf{A validity-closed molecular process.} Connectivity and the declared valence grammar are invariants of every committed event. Universal primitive operations are accelerated by verified subtree-Graft and coordinated ring rewrites under the same executor.
    \item \textbf{Structured-noise transport and target-free inference.} A directly samplable carbon-tree prior, marginal-preserving flexible-size coupling, and certified source-to-target programs train a non-monotone generator sampled ancestrally without a target, trace, beam, or repair phase.
    \item \textbf{A falsifiable operator study.} We specify matched null/tree, primitive/Graft, and scalar/coordinated-ring ablations, together with all-state and operator-use audits that test whether the learned process genuinely revises molecules rather than collapsing to any-order growth.
\end{itemize}

To our knowledge, this is the first molecular generative framework to learn, by Generator Matching, a CTMC whose marked events are legal matches of executable chemistry-constrained graph rewrite rules. The qualified claim is the synthesis---not the invention of graph rewriting, chemical rules, CTMCs, insertion/deletion, or Generator Matching individually.

\begin{figure}[t]
\centering
\fbox{\begin{minipage}{0.95\linewidth}
\centering
\textbf{Training}\qquad
$G_0\sim\pzero,\;G_1\sim\pdata$
$\xrightarrow{\text{couple + compile}}$
$z=(G_0\xrightarrow{a_1}\cdots\xrightarrow{a_K}G_1)$

\vspace{5pt}
$z\xrightarrow{\text{sample }(t,k)}$
$(G_k,t,\;q_t^z)$
$\xrightarrow{\text{Generator Matching}}$
$q_\theta(a\mid G_k,t)$

\vspace{5pt}
\textbf{Inference}\qquad
$\widetilde G_0\sim\pzero$
$\xrightarrow[\text{no endpoint; no compiler; no beam}]{\text{ancestral learned rewrite CTMC}}$
$\widetilde G_T\sim p_\theta$
\end{minipage}}
\caption{\textbf{Method overview.} The endpoint and rewrite program are latent training variables. Generator Matching learns their state-conditional marginal generator; inference repeatedly samples and executes legal marked rewrites.}
\label{fig:overview}
\end{figure}

\section{Background and Positioning}

\subsection{Generators and Generator Matching}

The infinitesimal generator characterizes a Markov process through its action on test functions. A deterministic flow has $(L_tf)(x)=v_t(x)^\top\nabla f(x)$, a diffusion adds a second-order term, and a pure jump process has
\begin{equation}
 (L_tf)(x)=\sum_{y\neq x}Q_t(x,y)[f(y)-f(x)].
\end{equation}
Generator Matching is a learning principle across these process classes \citep{holderrieth2025generator}: construct tractable endpoint-conditional generators, then regress their conditional expectation given the current state. It does not prescribe a molecular representation, allowed transitions, or source-to-target coupling. \rgm contributes those objects for rewrite systems; \method contributes their chemistry-specific implementation.

\subsection{Graph diffusion, flows, and flexible size}

DiGress and related discrete diffusion models corrupt node and edge categories and learn a denoising process \citep{vignac2023digress,austin2021d3pm,campbell2022ctdd}; DeFoG formulates graph generation by discrete flow matching \citep{qin2025defog}. Their intermediate tensors are discrete graphs but are not restricted to connected valence-admissible molecules. GrIDDD adds monotone node insertion and deletion to graph diffusion \citep{ninniri2025griddd}. Edit Flows constructs insertion, deletion, and substitution flows for variable-length objects \citep{havasi2025editflows}. Morph uses unbalanced optimal transport to change the cardinality of 3D geometric molecular graphs \citep{franke2026morph}. DDSBM learns CTMC bridges for graph transformation using independent node/edge modifications and iterative bridge fitting \citep{kim2025ddsbm}. These are close in flexibility or mathematical machinery. \method differs in the support and semantics of its generator: one jump is an executable, guarded chemical rewrite between complete molecules, and de novo inference does not condition on a paired endpoint.

Morph is the closest overlap in its use of Generator Matching, structural priors, and insertion/deletion during molecular generation. Accordingly, neither flexible cardinality, prior-to-data transport, nor Generator Matching by itself is a claimed novelty here. The distinction is that \rgm makes a legal rule match---rather than an independently parameterized node, edge, or coordinate update---the marked infinitesimal event, quotients syntactic aliases at the molecular-successor level, and samples entirely within the rewrite-defined chemical state space.

\subsection{Constraints, grammars, and stochastic rewriting}

ConStruct projects a discrete graph diffusion process onto selected structural constraints \citep{madeira2024construct}; CoCoGraph uses connectivity-preserving double-edge swaps inside fixed composition/degree fibers \citep{ruizbotella2026cocograph}. Molecular hypergraph grammars can make validity structural and improve data efficiency \citep{kajino2019mhg,guo2022grammar}; fragment and junction-tree models generate through a learned or predefined motif vocabulary \citep{jin2018jtvae}. Stochastic graph rewriting provides the closest operational precedent: qualitative rules define matches and application conditions, while contextual propensities induce a CTMC \citep{danos2015thermodynamic,behr2021rewriting}. \method imports this rule/rate separation, but replaces fixed or thermodynamic propensities with a time-inhomogeneous neural generator learned from endpoint-conditioned programs. Its primitive operators, deletion and correction semantics, and source-independent sampling also distinguish it from a fragment assembly ontology.

\begin{table}[t]
\centering
\caption{Conceptual comparison. ``Path support'' describes visible generative states, not endpoint validity. ``Size edits'' means changing cardinality within one trajectory, rather than drawing a size before generation.}
\label{tab:positioning}
\small
\begin{tabularx}{\textwidth}{lYYYY}
\toprule
Family & Learned infinitesimal object & Path support & Size edits & Typical endpoint information\\
\midrule
DiGress / DeFoG & node/edge transition rates & categorical graph tensors & no & none\\
DDSBM & node/edge bridge rates & categorical graph tensors & no & target distribution / bridge\\
GrIDDD / Edit Flows & insert, delete, and substitute rates & variable-cardinality objects & yes & none\\
Morph & unbalanced geometric transport & flexible-size 3D graphs & yes & none or condition\\
Molecular grammars & derivation policy & grammar-valid partial derivations & yes & none\\
Stochastic rewriting & rule-match propensity & rule-defined graph states & rule-dependent & none\\
\midrule
\method & legal rewrite-match rate & complete valid connected molecules & yes & none at inference\\
\bottomrule
\end{tabularx}
\end{table}

\section{Rewrite Generator Matching}

\subsection{Molecular states, legal match fibers, and quotient rates}

Let a padded heavy-atom graph be $G=(V,E,\ell,b)$, where $\ell_v$ contains element, formal charge, and implicit-hydrogen state and $b_{uv}$ is a discrete bond class. Given chemistry grammar $\Gamma$ and atom cap $N_{\max}$, define
\begin{equation}
\M=\{G:1\le |V|\le N_{\max},\;G\text{ is connected and valence-admissible under }\Gamma\}.
\label{eq:state-space}
\end{equation}
This is a representation-level domain: it does not assert synthesizability, low strain, or 3D feasibility.

A rule $r\in\Rset$ has a left pattern, a right pattern, and application conditions. A concrete mark $a=(r,m,\xi)$ consists of a match $m\in\Match_r(G)$ and a typed payload $\xi$ (for example a new atom state, bond order, or ring program). Its executor $\T_a$ is partial. The legal marked-action fiber is
\begin{equation}
\A(G)=\{(r,m,\xi):m\in\Match_r(G),\;\legal(G,r,m,\xi)=1\},
\quad \T_aG\in\M\;\;\forall a\in\A(G).
\label{eq:legal-fiber}
\end{equation}
A neural marked rate $q_\theta(a\mid G,t)\ge0$ defines
\begin{equation}
(L_{\theta,t}f)(G)=\sum_{a\in\A(G)}q_\theta(a\mid G,t)[f(\T_aG)-f(G)].
\label{eq:marked-generator}
\end{equation}
The marked representation is not unique: graph automorphisms, Kekul\'e alternatives, or distinct macro lowerings can name the same successor. Let $a\sim_G a'$ iff $\T_aG$ and $\T_{a'}G$ are the same canonical chemical state. The quotient successor fiber and its rate are
\begin{align}
\Sset(G)&=\{H\in\M:\exists a\in\A(G),\;H\simeq\T_aG\},\\
Q_{\theta,t}(G,H)&=\sum_{a\in\A(G):\;\T_aG\simeq H}q_\theta(a\mid G,t).
\label{eq:quotient-rate}
\end{align}
Equation~\eqref{eq:quotient-rate} is essential: a model may parameterize named matches, but probability evolves on chemical states through the sum of all aliases.

\subsection{Latent rewrite programs as conditional probability paths}

Sample a source--target pair from a coupling $\gamma$ with marginals $\pzero$ and $\pdata$. A randomized compiler samples a certified program
\begin{equation}
Z=(G_0,a_1,G_1^Z,\ldots,a_K,G_K^Z),\qquad
G_{k+1}^Z=\T_{a_{k+1}}G_k^Z\in\M,quad G_K^Z\simeq G_1.
\label{eq:program}
\end{equation}
The superscript distinguishes program states from the paired endpoint. The compiler distribution is part of the training coupling; multiple programs or source draws may be sampled for one endpoint.

Let $\alpha:[0,1]\to[0,1]$ be increasing with $\alpha(0)=0$ and $\alpha(1)=1$. Conditional on a length-$K$ program, define progress $K_t\sim\operatorname{Binomial}(K,\alpha(t))$ and $X_t=G_{K_t}^Z$. The pure-birth realization fires only the next instruction at latent progress $k<K$, with marked rate
\begin{equation}
r_t^Z(a\mid k)=
\one\{a=a_{k+1}\}(K-k)\frac{\alpha'(t)}{1-\alpha(t)}.
\label{eq:conditional-rate}
\end{equation}
For $\alpha(t)=t$, operational time $s=-\log(1-t)$ removes the terminal singularity and gives rate $K-k$. Sampling $k$ directly from its binomial marginal avoids simulating the teacher chain during training. If a program revisits the same molecular state, $k$ remains part of the latent variable; the learned rate conditions only on $X_t$ and therefore averages over those occurrences.

The conditional successor rate is
\begin{equation}
R_t^Z(H\mid X_t,k)=\sum_{a:\T_aX_t\simeq H}r_t^Z(a\mid k),
\end{equation}
and the marginal generator to be learned is
\begin{equation}
Q_t(G,H)=\E\!\left[R_t^Z(H\mid X_t,K_t)\mid X_t=G\right].
\label{eq:marginal-generator}
\end{equation}
This conditional expectation is the step inherited from general Generator Matching. The rewrite state space, program construction, legal fibers, and successor quotient are specific to \rgm.

\subsection{Poisson--Bregman Generator Matching objective}

For $\phi(u)=u\log u-u$, the Bregman divergence between nonnegative rates gives, up to target-only constants,
\begin{equation}
\mathcal L_{\mathrm{RGM}}(\theta)=
\E_{t,Z,K_t}\left[
\sum_{H\in\Sset(X_t)}
Q_{\theta,t}(X_t,H)-
R_t^Z(H\mid X_t,K_t)\log Q_{\theta,t}(X_t,H)
\right].
\label{eq:gm-objective}
\end{equation}
The first term calibrates total firing mass; the second rewards the conditional teacher successor. Unlike a classification loss over the observed action, it is proper for rates. Training can oversample rare rule-family progress positions using a mixture proposal $\widetilde p(k\mid t,Z)$ and weight by $p(k\mid t,Z)/\widetilde p(k\mid t,Z)$; this changes variance, not the population objective.

\section{COMPOSE: A Molecular Rewrite Process}

\subsection{Primitive and derived operator systems}

\begin{table}[t]
\centering
\caption{Rewrite families. A derived rule is admitted only if its macro transaction and verified lowering satisfy the same executor contract.}
\label{tab:operators}
\small
\begin{tabularx}{\textwidth}{lXX}
\toprule
Family & Semantic effect & Generative role\\
\midrule
Grow / shrink & attach a typed atom; delete a removable atom & variable size with no disconnected state\\
Atom restate & jointly change element, charge, and implicit H & composition and local-valence correction\\
Bond edit & insert/delete a non-bridge bond or change bond class & ring repair and unsaturation\\
Subtree Graft & cut one bridge and reattach the detached component atomically & global topology revision without atom churn\\
Ring ear & insert/delete a path between anchors, including same-anchor cycles & coordinated cyclic, fused, bridged, and spiro topology\\
Ring restate & verified joint electronic restatement on a fixed cyclic system & aromatic and heteroatom coordination\\
\bottomrule
\end{tabularx}
\end{table}

The micro language contains leaf insertion/deletion, atom restatement, bond insertion/deletion, and bond reordering. It is expressive but can be statistically inefficient: a full ring requires coordinated local decisions, while tree topology may be changed only after destructive atom churn. Derived rules compress recurring valid programs into one semantic mark. They do not create a fragment vocabulary: payloads factorize over topology, anchors, atom/bond states, and span, and the micro rules retain support for uncommon structures.

\paragraph{Subtree Graft.}
Choose a bridge $(u,v)$, conceptually remove it, choose endpoints in the two induced components, and add a new cross-component bond in the same atomic transaction. The old bridge cannot be exposed as a visible disconnected state. Graft preserves atom count and the full detached component, providing direct supervision for correction of random source-tree wiring.

\paragraph{Coordinated rings.}
Root-cycle and attached-cycle rules introduce a closed cyclic block; an open-ear rule inserts a path between two existing anchors; a same-anchor ear introduces a spiro petal; and a ring-system restatement coordinates common bond-order/heteroatom patterns. Open-ear decomposition supplies a compact topological substrate for monocyclic, fused, and bridged biconnected blocks, while same-anchor composition covers spiro articulation. Aromaticity is represented by executable bond/atom states, not an unconstrained post-hoc label. All ring sizes remain structurally possible subject to $N_{\max}$ and grammar guards; corpus frequencies belong in learned rates, not hard support.

\subsection{A directly samplable carbon-tree prior}

The principal de novo prior first samples a heavy-atom count $N\sim\mu$, then a degree-bounded tree by a capped Pr\"ufer construction, and realizes it as a neutral, single-bonded carbon graph. Thus $G_0\sim\pzero$ is connected, valence-admissible, and sampled without a target. The random tree is structured noise, not a retained scaffold; every atom may be retyped, moved by Graft, deleted, or replaced.

Coupling source and target sizes can shorten teacher paths without leaking an endpoint at inference. In the size-matched experiment, $N_0=N_1$ and $N_1\sim\mu$; hence the unconditional source-size marginal is still $\mu$. Our first flexible-size compiler instead draws $N_0\sim\mu$ independently of $N_1$, which exactly matches the ancestral prior and directly supervises both grow and shrink, but can produce longer programs. A lower-variance alternative is a size-jitter kernel $J(N_0\mid N_1)$ satisfying stationarity $\sum_n\mu(n)J(m\mid n)=\mu(m)$; a reversible Metropolis kernel on atom counts is one construction. The source wiring remains independent given $N_0$. We separately evaluate a null prior, an independently sized tree, a size-matched tree, and stationary size-jitter so that topology and cardinality effects are not conflated.

\subsection{Certified program compilation}

For equal-sized source trees, the compiler aligns temporary slots, selects a target spanning tree that maximizes overlap with the source, and uses Graft edge exchanges to reach that scaffold. It then restates atom states and retained bond classes and installs omitted target chords through ring-ear rewrites. Size-mismatched programs first delete or insert atoms under connectivity guards. Every emitted action is executed immediately, every intermediate is validated, and the final canonical graph must equal the target.

This offline planning problem is distinct from sampling. The compiler may use deterministic tie-breaking, bounded search, or a constructive fallback because it sees both endpoints. The learned sampler sees none of these. Calling inference ``ancestral'' means that each next event is sampled from the current learned hazard and then becomes part of the conditioning state; it does not mean fixed-order autoregression.

\subsection{Factorized neural rates and exact support masks}

Enumerating and encoding every successor separately is prohibitive. \method encodes the molecular graph once and factorizes a marked rate as
\begin{equation}
q_\theta(a\mid G,t)=\lambda_\theta(G,t)\,
p_\theta(r\mid G,t)\,
p_\theta(m,\xi\mid r,G,t),
\label{eq:factorization}
\end{equation}
where both categorical factors are normalized after analytic legality masks. A message-passing encoder provides graph, node, pair, and ring-context features. The total-hazard head learns how fast events occur; family and operand heads learn which legal event fires. Only the sampled mark is instantiated and passed to the full executor. An exhaustive canonical-successor enumerator remains a small-system oracle for mask equivalence, alias aggregation, and sampling tests.

\subsection{Target-free ancestral CTMC sampling}

Starting at $G\sim\pzero$, the sampler evaluates $\lambda_\theta(G,t)$, draws a waiting time, samples a masked family and payload, commits the rewrite, and recomputes rates from the successor. Rates are piecewise frozen between time-grid boundaries. There is no target, path compiler, candidate bank, beam, or terminal repair. Algorithm~\ref{alg:sampling} is detailed in \cref{app:sampler}.

\begin{figure}[t]
\centering
\fbox{\begin{minipage}{0.94\linewidth}
\textbf{Ancestral rewrite sampling.} Sample $G\sim\pzero$ and set operational time $s=0$. While $s<S$ and the event budget is not exhausted: (i) predict total legal hazard and masked marked probabilities; (ii) draw an exponential waiting time; (iii) if it crosses a frozen-rate boundary, advance to that boundary; otherwise sample one legal marked rewrite, execute $G\leftarrow\T_aG$, and advance by the waiting time. Return the current complete molecule.
\end{minipage}}
\caption{Target-free inference. A returned state is meaningful at every stopping time because the process never leaves $\M$.}
\label{alg:sampling}
\end{figure}

\section{Theoretical Properties}

We separate support guarantees, which follow from the rewrite system, from distributional correctness, which additionally requires exact Generator Matching.

\begin{assumption}[Finite closed rewrite system]\label{ass:finite}
$N_{\max}$ and all type vocabularies are finite; every committed action satisfies \eqref{eq:legal-fiber}; and the total rate is locally bounded on the finite time horizon.
\end{assumption}

\begin{assumption}[Conditional path consistency]\label{ass:path}
The program sampler has endpoint marginals $G_0\sim\pzero$ and $G_K^Z\sim\pdata$, and every sampled program satisfies \eqref{eq:program}.
\end{assumption}

\begin{proposition}[Pathwise molecular closure]\label{prop:closure}
Under \cref{ass:finite}, if $G_0\in\M$, every state committed by the teacher or learned CTMC lies in $\M$ almost surely.
\end{proposition}

\begin{proposition}[Connected topological reachability]\label{prop:reachability}
Ignoring chemical labels, any two finite connected simple graphs can be connected by a finite sequence of leaf insertions/deletions and non-bridge edge insertions/deletions, with every intermediate connected. For equal-order trees, subtree Grafts suffice using at most $(n-1)-|E(T_0)\cap E(T_1)|$ edge exchanges under a fixed vertex alignment.
\end{proposition}

\begin{remark}
The first statement is topological, not a universal chemistry theorem. In \method, every proposed step must also satisfy the valence grammar. Dataset-pair reachability is therefore certified constructively by compiler replay rather than inferred from \cref{prop:reachability} alone.
\end{remark}

\begin{proposition}[Population RGM minimizer]\label{prop:gm}
For every occupied $(G,t)$ and successor $H$, the unrestricted nonnegative-rate minimizer of \eqref{eq:gm-objective} is the marginal quotient rate $Q_t(G,H)$ in \eqref{eq:marginal-generator}.
\end{proposition}

\begin{corollary}[Endpoint law under an exact model]\label{cor:endpoint}
Under \cref{ass:finite,ass:path}, if $Q_{\theta,t}=Q_t$ on occupied states and the learned process starts at $\pzero$, its time marginals match the program-mixture marginals; in particular its terminal law is $\pdata$.
\end{corollary}

The proofs are in \cref{app:proofs}. These claims do not imply that a finite network attains the population minimizer, that teacher programs are optimal, or that 2D validity implies synthetic accessibility.

\section{Experiments}

\subsection{Questions, hypotheses, and protocol}

The evaluation is organized around five falsifiable questions.
\begin{enumerate}[leftmargin=*]
    \item[Q1.] \textbf{Distribution learning:} does \method match unconditional molecular distributions competitively under a matched GuacaMol protocol?
    \item[Q2.] \textbf{Support:} do all accepted intermediate states remain valid and connected, and are there any executor/mask discrepancies?
    \item[Q3.] \textbf{Topology:} does Graft shorten random-tree transport, prevent delete-to-one collapse, and improve endpoint topology statistics?
    \item[Q4.] \textbf{Rings:} do coordinated ring rules improve aromatic, saturated, fused, bridged, spiro, and ring-size fidelity beyond scalar bond edits?
    \item[Q5.] \textbf{Non-monotone use:} under flexible-size coupling, does the learned generator actually use grow, shrink, retype, bond revision, and Graft across time?
\end{enumerate}

\paragraph{Data and preprocessing.}
The principal benchmark is the GuacaMol distribution-learning corpus \citep{brown2019guacamol}. We canonicalize with a frozen RDKit version, declare the atom/charge/bond vocabulary and $N_{\max}$, and hash the accepted split. Final train/validation/test counts are $\XXX/\XXX/\XXX$. All exclusions and typed-ring catalog coverage are reported by split; validation and test molecules never construct proposal vocabularies or source-size statistics.

\paragraph{Metrics.}
Endpoint metrics are validity, uniqueness, novelty, Fr\'echet ChemNet Distance (FCD) \citep{preuer2018fcd}, GuacaMol KL, scaffold diversity \citep{bemis1996frameworks}, size and element total variation, and ring-topology distances. FCD uses exactly \XXX generated and \XXX reference molecules, with mean and covariance components and multi-seed uncertainty. Path metrics count every committed state, zero-hazard stalls, event-budget exhaustion, per-family events, source-edge survival, and invalid attempts. Because masked actions should be impossible, any all-step failure is reported as a software violation rather than averaged away.

\paragraph{Baselines and fairness.}
We compare with DiGress, DeFoG, GrIDDD, Morph, CoCoGraph, and grammar/fragment systems only where an official artifact can be evaluated under compatible chemistry and splits. Rerun and paper-reported values are labeled separately. DDSBM is evaluated as a graph-transformation neighbor, not presented as an unconditional de novo baseline. Parameter count, preprocessing, accelerator hours, sampling evaluations, and sanitization policy are reported.

\subsection{Main results: confirmatory evaluation pending}

\begin{table}[t]
\centering
\caption{Unconditional GuacaMol distribution learning. \textbf{Every entry is intentionally XXX until a frozen matched evaluation is complete}; no developmental value is promoted into this table.}
\label{tab:main}
\small
\begin{tabular*}{\textwidth}{@{\extracolsep{\fill}}lcccccc}
\toprule
Method & Valid $\uparrow$ & Unique $\uparrow$ & Novel $\uparrow$ & FCD $\downarrow$ & KL $\uparrow$ & Scaffold div. $\uparrow$\\
\midrule
DiGress & XXX & XXX & XXX & XXX & XXX & XXX\\
DeFoG & XXX & XXX & XXX & XXX & XXX & XXX\\
CoCoGraph & XXX & XXX & XXX & XXX & XXX & XXX\\
GrIDDD & XXX & XXX & XXX & XXX & XXX & XXX\\
Morph & XXX & XXX & XXX & XXX & XXX & XXX\\
\midrule
\method & \ComposeValidity & \ComposeUniqueness & \ComposeNovelty & \ComposeFCD & \ComposeKL & XXX\\
\bottomrule
\end{tabular*}
\end{table}

\begin{table}[t]
\centering
\caption{Path and structural fidelity. All-step metrics audit every accepted intermediate, not only endpoints. Values are pending confirmatory evaluation.}
\label{tab:path}
\small
\begin{tabular*}{\textwidth}{@{\extracolsep{\fill}}lccccccc}
\toprule
Variant & Step valid & Step conn. & Size TV & Cycle TV & Ring TV & Budget & Events\\
\midrule
Null prior & XXX & XXX & XXX & XXX & XXX & XXX & XXX\\
Independent tree & XXX & XXX & XXX & XXX & XXX & XXX & XXX\\
Size-matched Graft & XXX & XXX & XXX & XXX & XXX & XXX & XXX\\
Flexible-size Graft & \ComposeTrajectoryValidity & \ComposeTrajectoryConnectivity & \ComposeSizeTV & \ComposeCycleRankTV & \ComposeRingSizeTV & XXX & \ComposeMeanEvents\\
\bottomrule
\end{tabular*}
\end{table}

\subsection{Ablations and process diagnostics}

\begin{table}[t]
\centering
\caption{Matched ablations. Path length is measured on compiled teachers; FCD and ring TV are measured from target-free learned samples. All values remain XXX pending the frozen study.}
\label{tab:ablations}
\small
\begin{tabular}{lccc}
\toprule
Variant & FCD $\downarrow$ & Mean teacher length $\downarrow$ & Ring-size TV $\downarrow$\\
\midrule
Null source & \NullPriorFCD & XXX & XXX\\
Primitive tree (delete/regrow) & \PrimitiveTreeFCD & XXX & XXX\\
Tree + Graft & \GraftTreeFCD & \ComposePathLength & XXX\\
No coordinated ring macros & \NoMacroFCD & XXX & XXX\\
No successor alias aggregation & \NoAliasFCD & XXX & XXX\\
No importance-correct stratification & \NoStratificationFCD & XXX & XXX\\
\bottomrule
\end{tabular}
\end{table}

Reachability alone does not establish learned capability. We therefore report event fractions for Graft (\ComposeGraftFraction), grow (\ComposeGrowFraction), shrink (\ComposeShrinkFraction), retype (\ComposeRetypeFraction), and ring events (\ComposeRingFraction), stratified by model time and source/endpoint size. We additionally measure whether source edges survive, whether paths visit the one-atom state, and whether samples have the covariance collapse that aggregate FCD can obscure.

\subsection{Developmental evidence is not a benchmark result}

Current engineering evidence establishes executability, not final quality. In a developmental size-matched Graft run, the compiler produced 99,738 supported programs for 49,869 of 50,000 targets; a 3,000-update checkpoint reduced held-out RGM loss from 84.09 to 28.14 and reached 90.2\% teacher family accuracy. Separately, an early 1,000-sample gate produced 100\% valid/connected accepted rollouts on its restricted chemistry. These values diagnose the implementation and motivate the confirmatory experiment. They are not directly comparable to full GuacaMol baselines, are not inserted into Tables~\ref{tab:main}--\ref{tab:ablations}, and will be replaced by frozen multi-seed results.

\section{Discussion and Limitations}

The intended contribution is a new division of labor. Rewriting defines a valid state graph and executable event semantics; Generator Matching learns how probability should flow on that graph. The practical promise is strongest where intermediate molecules are consumed: online guidance can score every state, scaffold constraints can remove actions from the legal fiber, and source-conditioned editing can reuse the same correction operators. These are capabilities of the substrate, not evidence that every conditional task is solved. We therefore treat editing and property guidance as secondary evaluations rather than claiming them ``for free.''

The current state is a 2D heavy-atom graph with implicit hydrogens. Stereochemistry, conformers, radicals, metals, and multi-center bonding require extensions. Declared valence and RDKit sanitization are not guarantees of synthesizability. A finite typed ring catalog can under-cover rare teacher programs; the factorized ring decoder must be evaluated against that failure mode. The compiler introduces inductive bias and may concentrate supervision on one of many valid paths. A carbon-tree prior is useful only if it improves learned transport rather than merely shortening teachers. Finally, hard support guarantees do not imply low FCD, likelihood optimality, rapid mixing, or chemical usefulness.

\section{Conclusion}

\method turns stochastic molecular graph rewriting from a hand-designed simulator or grammar into the transition substrate of a learned generative process. Endpoint-conditioned valid programs provide exact jump-rate targets; Generator Matching removes the endpoint and program; guarded execution preserves the declared chemical state space at every committed event. Subtree Grafts and coordinated ring rules make the universal operator story practical without replacing it with fragment assembly. The resulting model is a target-free ancestral generator and a foundation for auditable molecular editing---subject to the empirical tests, marked XXX, that determine whether its distributional quality matches its structural guarantees.

\bibliography{references}
\bibliographystyle{iclr2026_conference}

\appendix

\section{Operator Semantics and Contracts}\label{app:operators}

For each $r$, $\Match_r(G)$ is the finite set of pattern monomorphisms satisfying structural application conditions. The payload $\xi$ contains new atom/bond states or a macro program. Rule, match, and payload define a partial map $\T_{r,m,\xi}:\M\rightharpoonup\M$. The implementation follows a propose--validate--commit discipline: construct a candidate copy, update implicit hydrogens and bond states jointly, run local valence and global sanitization/connectivity checks, then commit or expose no transition.

\begin{table}[h]
\centering
\caption{Minimum application conditions. The implementation may impose additional grammar-specific checks.}
\small
\begin{tabularx}{\textwidth}{lX X}
\toprule
Rule & Structural condition & Chemical/transaction condition\\
\midrule
Atom insert & empty slot; occupied anchor & allowed atom/bond; endpoint valences feasible\\
Atom delete & removable atom; remainder connected & neighbor H compensation; sanitized successor\\
Atom restate & occupied slot & joint element/charge/H state allowed\\
Bond insert & absent edge; endpoints occupied & allowed order; valence feasible; nonidentity\\
Bond delete & existing non-bridge edge & endpoint H compensation; sanitized successor\\
Bond reorder & existing edge & joint order/H update feasible\\
Subtree Graft & old edge is a bridge; new endpoints cross its cut & atomic cut+attach; all endpoint valences feasible\\
Ring ear insert & valid anchors and fresh internal slots & joint typed path closes; full lowering replays\\
Ring ear delete & matched removable ear; remainder connected & exact inverse payload; full lowering replays\\
Ring restate & fixed matched cyclic system & coordinated atom/bond state passes executor\\
\bottomrule
\end{tabularx}
\end{table}

\paragraph{Why macros remain rewrite events.}
A macro has one chemically interpretable successor and one learned rate. Its lowering is a proof-carrying implementation artifact: each micro state must also lie in $\M$, but inference exposes the macro as one marked CTMC event. This avoids treating arbitrary edit bundles as simultaneous jumps. Unrelated micro events may conflict on valence, fail to commute, and create an exponential bundle space; only verified compound semantics are promoted.

\paragraph{Ring topology coverage.}
Every 2-connected simple graph admits an open-ear decomposition starting from a cycle. Consequently, root/attached cycles plus distinct-anchor ears span biconnected cyclic topology at the graph level. Spiro systems join blocks at an articulation and are represented by same-anchor cycles. Repeated composition covers fused, bridged, spiro, linked, and mixed polycycles. Chemical feasibility remains enforced per typed payload; the graph theorem does not certify strain or synthesis.

\section{Source Priors and Couplings}\label{app:prior}

\subsection{Degree-bounded carbon-tree sampler}

Given size law $\mu$ on $\{1,\ldots,N_{\max}\}$:
\begin{enumerate}[leftmargin=*]
    \item sample $N\sim\mu$;
    \item sample a Pr\"ufer sequence with rejection or capped construction so every degree is at most four;
    \item decode the tree, assign neutral carbon atom states and single bonds, and derive implicit hydrogens;
    \item assert connectivity, valence validity, and the declared slot count.
\end{enumerate}
The source sampler consumes only $\mu$ and an RNG seed. For a training pair, the tree draw is independent of the endpoint conditional on the selected source size.

\subsection{Stationary size jitter}

Let $J$ be a Markov kernel on sizes with stationary law $\mu$. Drawing $N_1\sim\mu$ then $N_0\sim J(\cdot\mid N_1)$ preserves the source marginal:
\begin{equation}
\Pr(N_0=m)=\sum_n\mu(n)J(m\mid n)=\mu(m).
\end{equation}
A local proposal $n\to n\pm1$ with a Metropolis--Hastings accept/reject correction is a simple reversible choice. The jitter scale controls teacher transport cost but must not be confused with the source law. At inference, $N_0$ is drawn directly from $\mu$ without sampling a target size.

\section{Program Compilation Algorithms}\label{app:compiler}

\subsection{Primitive connectivity-preserving transport}

For completeness, any source can be reduced topologically to one retained atom: delete non-tree edges to expose a spanning tree, then remove leaves. Reverse the construction for a target spanning tree, add chords, and apply typed restatements. The production compiler uses direct common structure and Grafts because the one-atom route is valid but statistically wasteful.

\subsection{Graft transport}

Given equal-order carbon tree $T_0$ and target $G_1$:
\begin{enumerate}[leftmargin=*]
    \item align temporary source slots with target slots and choose a spanning tree $T_1\subseteq G_1$ maximizing $|E(T_0)\cap E(T_1)|$;
    \item while $T\neq T_1$, select $e\in E(T_1)\setminus E(T)$; adding $e$ creates a unique cycle, from which select $f\in E(T)\setminus E(T_1)$; execute the corresponding atomic subtree Graft and validate;
    \item restate retained atom states and tree-bond classes in a guard-feasible order;
    \item install target chords with verified ring-ear or bond events;
    \item canonicalize and assert exact equality with $G_1$; replay the complete program independently.
\end{enumerate}
The graph-level edge-exchange loop performs at most $(n-1)-|E(T_0)\cap E(T_1)|$ iterations. Chemistry guards may reject a nominal exchange, so the implementation searches alternative cycle edges and uses a constructive fallback. A compilation failure is explicit; it is never converted to a truncated target.

\subsection{Ring catalog construction}\label{app:ring-catalog}

The current finite proposal catalog is mined only from training programs:
\begin{enumerate}[leftmargin=*]
    \item decompose each target into biconnected blocks and its block--cut tree;
    \item compile cyclic blocks into root/attached cycles, open ears, same-anchor cycles, and electronic restatements;
    \item canonicalize marks modulo slot permutation and chemically equivalent aliases;
    \item count marks and retain a declared training-only vocabulary or factorized fields;
    \item audit full-program coverage independently on train, validation, and test.
\end{enumerate}
Frequency initializes or stratifies learning but does not define structural legality. The final factorized decoder predicts family, anchors, span, and typed atom/bond payload so that uncommon combinations need not occur as memorized fragments.

\subsection{Program-cache reproducibility}

Programs are stored as rewrite instructions plus byte-packed exact graph checkpoints rather than a dense graph at every progress position. A sampled state replays at most $c-1$ validated actions from the nearest checkpoint. CPU compilation writes atomic, signature-validated shards with a manifest containing split hashes, compiler version, chemistry grammar, prior/coupling configuration, ring catalog hash, seeds, and checkpoint interval. A GPU training stage refuses to start from an incomplete cache. These choices change throughput and storage, not the mathematical path distribution.

\section{Proofs}\label{app:proofs}

\subsection{Proof of \cref{prop:closure}}

Induct on committed events. The base state is in $\M$. If $G_n\in\M$, the CTMC has support only on $a\in\A(G_n)$; by \eqref{eq:legal-fiber}, $G_{n+1}=\T_aG_n\in\M$. Local boundedness excludes explosion on the finite horizon, so only finitely many jumps occur almost surely. The conclusion applies equally to teacher programs and learned sampling because both invoke the same executor. \qed

\subsection{Proof of \cref{prop:reachability}}

For connected $G$, choose a spanning tree $T_G$. Delete every edge in $E(G)\setminus E(T_G)$; each is on a cycle at deletion and is therefore not a bridge, so connectivity is preserved. Repeatedly delete leaves until one vertex remains. From one vertex, reverse leaf insertions to construct a spanning tree $T_H$ of $H$, then insert edges in $E(H)\setminus E(T_H)$. Every intermediate is connected.

For aligned trees $T$ and $T_1$, take $e\in E(T_1)\setminus E(T)$. Adding $e$ creates a unique cycle. That cycle contains an edge $f\in E(T)\setminus E(T_1)$; otherwise $T_1$ would contain a cycle. Replacing $f$ by $e$ is one subtree Graft, preserves the tree property, and increases overlap with $T_1$ by one. Repetition terminates after at most $(n-1)-|E(T_0)\cap E(T_1)|$ exchanges. \qed

\subsection{Proof of \cref{prop:gm}}

Fix occupied $(G,t,H)$ and let $Y=R_t^Z(H\mid X_t,K_t)$ conditional on $X_t=G$. The part of \eqref{eq:gm-objective} depending on scalar prediction $u\ge0$ is
\begin{equation}
\ell(u)=u-\E[Y\mid X_t=G]\log u,
\end{equation}
with target-only terms omitted. For $u>0$, $\ell'(u)=1-\E[Y\mid X_t=G]/u$ and $\ell''(u)=\E[Y\mid X_t=G]/u^2\ge0$. If the mean is positive, the unique minimizer is its value; if zero, the infimum is attained at $u=0$ by continuity. Finite sums commute with conditional expectation, so the result applies after alias aggregation. \qed

\subsection{Proof of \cref{cor:endpoint}}

The program mixture induces marginal law $p_t$ satisfying the Kolmogorov forward equation with $Q_t$ from \eqref{eq:marginal-generator}; this is the standard conditional-to-marginal Generator Matching identity. Under \cref{ass:finite}, the finite-state forward equation has a unique solution. If $Q_{\theta,t}=Q_t$ and both laws begin at $\pzero$, their marginals agree for all $t$. By \cref{ass:path}, the program mixture terminates at $\pdata$, hence so does the exact learned process. \qed

\subsection{Importance-correct progress stratification}

Let $p(k\mid t,Z)$ be the binomial progress law and $h(k\mid Z)$ a family-balanced proposal. Sample from $\widetilde p=(1-\rho)p+\rho h$ and weight by $w=p/\widetilde p$. For any integrable per-progress loss $g$, $\E_{\widetilde p}[wg]=\sum_kp(k)g(k)=\E_p[g]$. Since $\widetilde p\ge(1-\rho)p$, support is retained for $\rho<1$.

\section{Model, Training, and Sampling Details}\label{app:model}

\subsection{Neural architecture}

The graph encoder embeds atom type, formal charge, implicit hydrogen count, bond class, model time, and ring/topology context. Message passing produces node and graph representations. A graph head predicts total hazard; a masked family head predicts the rule schema; atom, pair, and template heads predict operands and typed payloads. Hierarchical normalization prevents a family with many raw matches from receiving probability solely because of multiplicity. Exact final values are: hidden width \XXX, layers \XXX, parameters \ComposeParameters, batch size \XXX, optimizer \XXX, learning rate \XXX, time schedule \XXX, and training updates \XXX.

\subsection{Training estimator}

For each minibatch item:
\begin{enumerate}[leftmargin=*]
    \item draw or load $(G_0,G_1,Z)$ from the declared coupling and compiled cache;
    \item sample time $t$ and progress $k$, optionally from the importance-correct stratified proposal;
    \item reconstruct $X_t=G_k^Z$ from its nearest packed checkpoint;
    \item encode $X_t$ once, form analytic legality masks, and evaluate \eqref{eq:gm-objective};
    \item aggregate aliases for the teacher successor and update $\theta$.
\end{enumerate}
We log hazard, family/operand accuracy, illegal-mask disagreements, gradient norm, and held-out RGM loss. Model selection uses a declared validation criterion without test FCD.

\subsection{Exact ancestral sampler}\label{app:sampler}

Let $0=s_0<\cdots<s_J=S$ be the operational-time grid. Starting from $G\sim\pzero$, at grid interval $j$:
\begin{enumerate}[leftmargin=*]
    \item compute $\lambda=\lambda_\theta(G,s)$ and masked marked probabilities;
    \item if $\lambda=0$, advance to $s_{j+1}$ and log a stall;
    \item otherwise draw $\Delta\sim\mathrm{Exponential}(\lambda)$;
    \item if $s+\Delta\ge s_{j+1}$, set $s=s_{j+1}$ without an event;
    \item else sample $a$ hierarchically, execute $G\leftarrow\T_aG$, set $s\leftarrow s+\Delta$, and recompute all rates.
\end{enumerate}
The event budget is a numerical safeguard and its hit rate is reported. There is no rejection loop over invalid actions: analytic masks and the executor define zero support, and any disagreement is an auditable implementation error.

\subsection{Complexity}

Let $n\le N_{\max}$, $d$ be hidden width, $L$ message-passing layers, and $C_r(G)$ the tensorized candidate count for family $r$. One rate evaluation costs the graph encoder plus masked heads, approximately $O(L(n d^2+|E|d)+\sum_r C_r(G)d)$ for the dense implementation, rather than one graph encoding per canonical successor. Compilation is offline and program-dependent; Graft tree alignment is polynomial and each edge exchange is validated. Final measured preprocessing time, peak memory, train throughput, and sample throughput are \ComposeTrainHours, \ComposePeakMemory, \XXX, and \ComposeThroughput.

\section{Complete Experimental Protocol}\label{app:experiments}

\subsection{Frozen data contract}

We will release raw input hashes, canonicalization code, exact accepted SMILES, split indices, $N_{\max}$, element/charge/bond vocabularies, every exclusion count, and the ring-catalog audit. Test molecules do not influence the source-size law, templates, hyperparameters, or checkpoint selection.

\subsection{FCD contract}

FCD is computed with one frozen ChemNet implementation and environment. Invalid or null samples are not replaced. We report generated/reference counts, total FCD, additive mean and covariance terms, generated/reference covariance traces, multi-seed mean and standard deviation, and bootstrap intervals. We use at least \XXX independent samples per final seed; any 2,000-sample diagnostic is labeled preliminary.

\subsection{Ablation protocols}

\paragraph{Prior/coupling.}
Hold architecture, data, steps, and sampler fixed; compare null, independently sized carbon tree, size-matched tree, and stationary size-jitter. Report teacher length as well as learned endpoint metrics.

\paragraph{Graft.}
Compare primitive delete/regrow transport with Graft-enabled transport on identical sampled source--target pairs. Measure compilation success, path length, one-atom visits, source-edge survival, training throughput, event use, and endpoint FCD/topology.

\paragraph{Rings.}
Compare scalar bond construction, topology-only ears, typed ears, and typed ears plus ring-system restatement. Report per-class recall/TV for aromatic, saturated, fused, bridged, spiro, small rings (3--4), ordinary rings (5--7), and macrocycles under the dataset definition.

\paragraph{Objective and quotient.}
Ablate total-hazard learning, importance-correct stratification, and canonical alias aggregation. The last ablation tests whether symmetric syntactic multiplicity biases chemical successor rates.

\subsection{Trajectory audits}

For every sampled event we record time, rule family, operands, pre/post canonical keys, hazard, and validity/connectivity flags. Aggregates include:
\begin{itemize}[leftmargin=*]
    \item rule-family and inverse-pair utilization by time and size;
    \item source-edge survival, topology edit distance, and one-atom-state visitation;
    \item ring class before and after every cyclic event;
    \item aliases per successor and marked-versus-quotient calibration;
    \item event-budget exhaustion, zero-hazard stalls, and time-grid sensitivity;
    \item agreement of analytic masks/direct sampling with exhaustive fibers on small graphs.
\end{itemize}

\subsection{Additional pending results}

\begin{table}[h]
\centering
\caption{Ring taxonomy distribution distances; all values pending.}
\small
\begin{tabular}{lcccccc}
\toprule
Method & Aromatic & Saturated & Fused & Bridged & Spiro & Macrocycle\\
\midrule
DiGress & XXX & XXX & XXX & XXX & XXX & XXX\\
Morph & XXX & XXX & XXX & XXX & XXX & XXX\\
\method & XXX & XXX & XXX & XXX & XXX & XXX\\
\bottomrule
\end{tabular}
\end{table}

\begin{table}[h]
\centering
\caption{Compute and storage; preprocessing is included and all values are pending.}
\small
\begin{tabular}{lccccc}
\toprule
Method & Parameters & Compile h & Train h & Peak memory & Molecules/s\\
\midrule
\method & \ComposeParameters & XXX & \ComposeTrainHours & \ComposePeakMemory & \ComposeThroughput\\
\bottomrule
\end{tabular}
\end{table}

\subsection{Conditional editing protocol}

Conditional editing is secondary to unconditional distribution learning. We initialize from held-out molecules, optionally freeze a scaffold by removing incompatible actions from $\A(G)$, and evaluate property improvement, scaffold retention, edit count, diversity, and all-step validity. The final success, retention, and gain values are \ComposeConditionalSuccess, \ComposeScaffoldRetention, and \ComposePropertyGain. Reward fine-tuning or inference guidance must be reported separately from the unguided generator.

\section{Failure Modes, Broader Impact, and Reproducibility}

\paragraph{Failure modes.}
We stratify failures into compiler support, proposal-mask false negatives, executor discrepancies, zero-hazard stalls, event-budget exhaustion, source-memory artifacts, invalid ring electronics, over-fusion, macrocycle excess, and endpoint distribution mismatch. Chemical examples are selected by predeclared strata rather than only visual appeal.

\paragraph{Broader impact.}
The model can generate valid but unsafe, toxic, or controlled compounds. A distribution-learning score is not evidence of therapeutic value. Downstream property guidance requires access control, dual-use review, toxicity assessment, and wet-lab governance. Corpus bias and cheminformatics blind spots remain even when formal validity is guaranteed.

\paragraph{Artifact checklist.}
The final release will contain:
\begin{itemize}[leftmargin=*]
    \item code version, environment lockfile, and RDKit/ChemNet versions: \XXX;
    \item dataset, split, vocabulary, and canonicalization hashes: \XXX;
    \item all hyperparameters, seeds, source draws, and coupling settings: \XXX;
    \item compiled-program manifest, ring catalog, and recovery checkpoints: \XXX;
    \item generated SMILES, per-event trajectories, and raw metric outputs: \XXX;
    \item multi-seed summaries, confidence intervals, and failed-run disclosure: \XXX.
\end{itemize}

\section{Large Language Model Usage Statement}

Large language models were used for drafting and code assistance. All mathematical claims, citations, implementation details, and empirical values require author verification. No generated reference or result is retained without checking a primary source or frozen experiment artifact.

\end{document}
